\section{问题三：基于连续制氨调节的运行分析}

\subsection{模型建立}

问题三仍考虑 24 种风光组合出力场景，但将问题二中的离散开停机调度
推广为连续功率调节。设 $\alpha_{t,s,q}$ 为场景 $s$、目标日产量 $q$
下第 $t$ 个时段电氢氨装置的负荷率。根据题设，装置功率连续可调且下限为
10\%，因此有
\begin{equation}
  0.1 \leq \alpha_{t,s,q} \leq 1,
  \qquad t=1,2,\dots,24.
  \label{eq:p3-alpha-bound}
\end{equation}
与问题二的二元变量 $u_{t,s,q}\in\{0,1\}$ 相比，式~\eqref{eq:p3-alpha-bound}
允许装置在最低负荷和满负荷之间连续变化，从而以更细的功率粒度匹配风光出力。

扩容至 $72~\mathrm{t/d}$ 后，设备满负荷功率沿用问题二参数。连续调节下
各设备功率和产氨速率表示为
\begin{equation}
  P_{t,s,q}^{\mathrm{ALK}}=20\alpha_{t,s,q},\quad
  P_{t,s,q}^{\mathrm{PEM}}=20\alpha_{t,s,q},\quad
  P_{t,s,q}^{\mathrm{NH_3}}=1.5\alpha_{t,s,q},\quad
  m_{t,s,q}^{\mathrm{NH_3}}=3\alpha_{t,s,q}.
  \label{eq:p3-device-power}
\end{equation}
因此，目标日产量约束为
\begin{equation}
  \sum_{t=1}^{24}3\alpha_{t,s,q}\Delta t=q,
  \qquad \Delta t=1~\mathrm{h}.
  \label{eq:p3-production}
\end{equation}
由于每小时至少以 10\% 负荷运行，任一方案均具有 $7.2~\mathrm{t/d}$
的基础产量；其余产量由各时段负荷率的连续增量分配完成。

园区内部用电功率为
\begin{equation}
  P_{t,s,q}^{\mathrm{use}}
  =P_{t,s}^{\mathrm{load}}+41.5\alpha_{t,s,q}.
  \label{eq:p3-use-power}
\end{equation}
购电功率和上网功率继续沿用问题一的功率平衡式~\eqref{eq:power-balance}
以及正负缺额定义式~\eqref{eq:p1-grid-exchange}。同时，为明确同一小时不能
既购电又售电，引入模式变量 $z_{t,s,q}\in\{0,1\}$，并设置
\begin{equation}
  0\leq P_{t,s,q}^{\mathrm{buy}}\leq Mz_{t,s,q},\qquad
  0\leq P_{t,s,q}^{\mathrm{sell}}\leq M(1-z_{t,s,q}),
  \label{eq:p3-grid-exclusive}
\end{equation}
其中 $M$ 为足够大的功率上界。式~\eqref{eq:p3-grid-exclusive} 保证任一时段
内购电和售电至多发生一种。

优化目标仍为最小化日运行净成本：
\begin{equation}
  \min C_{s,q}^{\mathrm{day}},
  \label{eq:p3-objective}
\end{equation}
其中 $C_{s,q}^{\mathrm{day}}$ 的构成与问题二式~\eqref{eq:p2-day-cost}
一致，包括风电和光伏度电成本、设备运维成本、分时购电费用以及余电上网
收入抵扣。吨氨成本继续采用式~\eqref{eq:p2-unit-cost} 的口径计算。
绿电直连指标则沿用问题一式~\eqref{eq:indicators}，并按全满足、部分满足、
全不满足三类统计。

从求解角度看，连续负荷率使各小时的可调空间可在 $0.1$ 到 $1$ 之间连续分配。
当
\begin{equation}
  \alpha_{t,s}^{0}
  =\frac{P_{t,s}^{\mathrm{re}}-P_{t,s}^{\mathrm{load}}}{41.5}
  \label{eq:p3-zero-exchange}
\end{equation}
落在 $[0.1,1]$ 内时，该负荷率对应本时段既不购电也不上网的临界状态。
由于谷电价可能低于余电上网电价，简单地将各小时边际成本独立排序可能会破坏
负荷率的连续路径。本文因此采用带买卖电互斥约束的混合整数线性规划求解，
其目标函数、功率平衡和产量约束均为线性形式，可直接求得连续调节下的全局最优方案。

\subsection{模型求解与结果}

对 24 种风光组合场景和五个目标日产量分别求解连续调节模型，并按每个场景代表
15 天折算全年。表~\ref{tab:p3-annual-result} 给出了固定日产量口径下的全年结果。
与问题二一致，$72~\mathrm{t/d}$ 方案因式~\eqref{eq:p3-production} 要求
$\sum_t\alpha_t=24$，各时段均为满负荷运行，因此与问题二结果相同。
在其他产量下，连续调节可在 24 小时内重新分配负荷率，使生产功率更贴近风光出力。

\begin{table}[H]
  \centering
  \caption{问题三 24 场景全年折算下不同日产量的运行结果}
  \label{tab:p3-annual-result}
  \begin{tabular}{ccccc}
    \toprule
    日产量 & 吨氨成本 & 网购电量 & 上网电量 & 指标天数统计 \\
    (t/d) & (元/t) & (MWh) & (MWh) & 全满足/部分/不满足 \\
    \midrule
    72 & 7228.97 & 221146.80 & 12078.72 & 270/90/0 \\
    63 & 6466.04 & 176326.80 & 12078.72 & 300/60/0 \\
    54 & 5725.68 & 137231.50 & 17803.42 & 240/120/0 \\
    45 & 5082.84 & 109563.64 & 34955.56 & 240/120/0 \\
    36 & 4302.07 & 91060.62 & 61272.54 & 165/195/0 \\
    \bottomrule
  \end{tabular}
\end{table}

从表~\ref{tab:p3-annual-result} 可见，若仅以吨氨成本最低为准，24 个场景下的
推荐方案仍倾向于选择 $36~\mathrm{t/d}$。该方案全年折算制氨量为 12960~t，
总净成本为 55\,754\,817.22~元，吨氨成本为 \textbf{4302.07~元/t}。
其新能源自发自用比例为 $64.24\%$、宽口径总用电绿电比例为 $41.95\%$，
均满足阈值；但新能源上网比例为 $35.76\%$，仍高于 $20\%$ 的要求。
这说明连续调节改善了新能源消纳，但在低日产量下，受总产量需求限制，
仍无法完全吸收高风光场景中的富余电量。

图~\ref{fig:p3-dispatch-example} 给出了代表性场景 $W4\_P1$、日产量
$36~\mathrm{t/d}$ 下的连续调度结果。可以看出，负荷率并非简单取 10\%
或 100\%，而是在不同小时之间连续变化：风光出力较高或购电成本较低时段提高负荷，
风光不足且购电成本较高时段降低负荷，从而同时压低购电量和上网量。

\begin{figure}[H]
  \centering
  \includegraphics[width=0.88\textwidth]{p3_continuous_schedule.pdf}
  \caption{问题三代表场景下连续负荷率调度与电网交换功率}
  \label{fig:p3-dispatch-example}
\end{figure}

问题三与问题二的全年吨氨成本对比如图~\ref{fig:p3-cost-compare} 所示。
除 $72~\mathrm{t/d}$ 因两种模型均为全天满负荷而结果相同外，其余产量下
连续调节均降低了吨氨成本。其中 $36~\mathrm{t/d}$ 方案的吨氨成本由
4586.42~元/t 降至 4302.07~元/t，下降 284.35~元/t；$45~\mathrm{t/d}$
方案下降 281.55~元/t。成本降低的主要原因在于，连续负荷可以利用离散开停机
难以覆盖的时段，使部分原本上网的新能源转化为内部制氨用电，并减少后续高价
时段的网购电量。

\begin{figure}[H]
  \centering
  \includegraphics[width=0.82\textwidth]{p3_p2_cost_comparison.pdf}
  \caption{问题二与问题三不同日产量下全年吨氨成本对比}
  \label{fig:p3-cost-compare}
\end{figure}

以成本最低的 $36~\mathrm{t/d}$ 方案为例，问题三相较问题二的全年网购电量
和余电上网量均减少 $1.62\times 10^4~\mathrm{MWh}$。对应三项绿电指标变化如图
~\ref{fig:p3-indicator-compare} 所示：新能源自发自用比例提高 9.46 个百分点，
宽口径总用电绿电比例提高 8.26 个百分点，上网比例降低 9.46 个百分点。
这表明，连续调节的作用并非改变总风光资源，而是改善风光出力与制氨负荷的
时序匹配关系。

\begin{figure}[H]
  \centering
  \includegraphics[width=0.78\textwidth]{p3_p2_indicator_comparison.pdf}
  \caption{问题二与问题三成本最低方案的绿电直连指标对比}
  \label{fig:p3-indicator-compare}
\end{figure}

综合来看，连续制氨调节相较离散开停机具有更好的消纳能力和经济性表现：
一方面，低负荷运行下限保证装置全天维持基本生产；另一方面，连续负荷率又允许
模型在新能源富余时段逐步提高功率，而不是被迫在停机和满负荷之间跳变。
因此，问题三模型在大多数产量档位下同时降低了购电量、上网量和吨氨成本。
不过，若要求三项绿电直连指标全部满足，单纯追求最低吨氨成本仍不充分；
例如 $36~\mathrm{t/d}$ 方案虽然成本最低，但上网比例仍超限。若优先考虑指标合规，
$54~\mathrm{t/d}$ 可作为更均衡的方案，其全年吨氨成本为 5725.68~元/t，
自发自用比例为 $89.61\%$、宽口径总用电绿电比例为 $49.76\%$、上网比例为
$10.39\%$，三项指标均满足要求。
